% ============================================================
%  Physical-Contact (Touch) Detection — Detailed Investigation Report
%  Smart Thermal System for Patient Safety Monitoring
%  Waveshare 26984 (80x62) Thermal Sensor Dataset
%  Afeka College of Engineering, Tel-Aviv
% ============================================================
\documentclass[11pt, a4paper]{article}

\usepackage[a4paper, margin=2.4cm]{geometry}
\usepackage[T1]{fontenc}
\usepackage{lmodern}
\usepackage{booktabs}
\usepackage{array}
\usepackage{multirow}
\usepackage{amsmath}
\usepackage{amssymb}
\usepackage[table]{xcolor}
\usepackage{graphicx}
\usepackage{hyperref}
\usepackage{parskip}
\usepackage{enumitem}
\usepackage{float}

\providecommand{\arraybackslash}{\let\\\tabularnewline}
\makeatletter
\providecommand\do@row@strut{}
\let\do@row@strut\relax
\makeatother

\definecolor{tpgreen}{RGB}{0,128,0}
\definecolor{tnblue}{RGB}{31,73,125}
\definecolor{fpred}{RGB}{192,0,0}
\definecolor{fnyellow}{RGB}{150,100,0}
\definecolor{headerblue}{RGB}{31,73,125}
\definecolor{cellgreenlight}{RGB}{198,239,206}
\definecolor{cellbluelight}{RGB}{221,235,247}
\definecolor{cellredlight}{RGB}{255,199,206}
\definecolor{cellyellowlight}{RGB}{255,235,156}
\definecolor{rowhi}{RGB}{223,240,216}
\hypersetup{colorlinks=true, linkcolor=headerblue, urlcolor=headerblue}

% \confmat{posLabel}{negLabel}{TP}{FN}{FP}{TN}
\newcommand{\confmat}[6]{%
    \vspace{2pt}\begin{center}\renewcommand{\arraystretch}{1.5}%
    \begin{tabular}{lcc}
        \hline\rule{0pt}{2.2ex}%
        & \textbf{Pred:\ #1} & \textbf{Pred:\ #2} \\[1pt]\hline\rule{0pt}{2.4ex}%
        \textbf{Truth:\ #1}
            & \cellcolor{cellgreenlight}\textcolor{tpgreen}{\textbf{TP}}\,$=$\,#3
            & \cellcolor{cellyellowlight}\textcolor{fnyellow}{\textbf{FN}}\,$=$\,#4 \\[3pt]
        \textbf{Truth:\ #2}
            & \cellcolor{cellredlight}\textcolor{fpred}{\textbf{FP}}\,$=$\,#5
            & \cellcolor{cellbluelight}\textcolor{tnblue}{\textbf{TN}}\,$=$\,#6 \\[1pt]\hline
    \end{tabular}\end{center}\vspace{2pt}}

\title{%
    {\large Afeka College of Engineering, Tel-Aviv}\\[0.3em]
    {\large Smart Thermal System for Patient Safety Monitoring}\\[0.7em]
    {\LARGE\bfseries Physical-Contact (Touch) Detection}\\[0.2em]
    {\large Detailed Investigation: Methods, Reasoning, and Results}\\[0.3em]
    {\normalsize Waveshare 26984 \quad 80$\times$62\,px \quad 3 synchronised views}}
\author{Guy Chen \and Yaniv Blau \and Roy Lieberman\\[0.2em]
    \small Supervisor: Or Zilberberg \qquad Advisor: Dr.\@ Oshrit Hoffer}
\date{June 2026}

\begin{document}
\maketitle
\tableofcontents
\newpage

% ============================================================
\section{Problem, data, and evaluation}
% ============================================================

\subsection{Task}
The touch module detects \emph{physical contact between patients} (covering the
``physical contact'' and ``violent/sexual behaviour'' scenarios) from three
synchronised MLX-class thermal cameras --- here the Waveshare 26984 at
80$\times$62\,px, 8\,Hz. No optical imagery is used, preserving privacy. A frame
is positive if two or more people are in physical contact in the scene.

\subsection{Labels and the central constraint}
Contact is annotated per frame in \texttt{contact\_labels.csv}. The decisive
fact for everything below is its scarcity:
\begin{itemize}[leftmargin=1.4em, itemsep=1pt]
\item \textbf{188 positive frames} in the entire dataset, concentrated in only
      \textbf{four scenes}: \texttt{2ppl\_fight}~(50), \texttt{3pp\_surprise}~(66),
      \texttt{3pplhedroncolider}~(53), \texttt{2ppl\_hug}~(19).
\item All other scenes are contact-negative and act as \textbf{hard negatives}:
      single people, and crucially people who are \emph{near but not touching}
      (\texttt{2ppl\_walk}, \texttt{3ppl\_dance}, \texttt{the\_more\_the\_merrier}).
\end{itemize}

\subsection{Evaluation protocol}
Every method is scored on a common per-session \emph{timeline} split (first 60\%
of each scene $\to$ train, next 15\% $\to$ validation, last 25\% $\to$ test),
giving a fixed \textbf{test set of 695 frames (41 positive, 654 negative)}.
Timeline (not random) splitting preserves temporal contiguity, required by the
temporal models. Metrics: accuracy, precision, recall, F1, and false-alarm rate
(FP/(FP+TN)). \textbf{F1} is primary; \textbf{recall} is the safety priority and
\textbf{false-alarm rate} guards against alert fatigue (itself a safety risk in a
ward). §\ref{sec:cv} revisits this single split with cross-validation.

\subsection{Result at a glance}
Table~\ref{tab:arc} previews the full arc; each row is detailed in its section.

\begin{table}[H]\centering
\caption{Progression of touch detectors on the common 695-frame test split.}
\label{tab:arc}
\renewcommand{\arraystretch}{1.2}
\begin{tabular}{l l c}
\toprule
\textbf{Method} & \textbf{Core idea} & \textbf{F1} \\
\midrule
Geometric multi-view (\S\ref{sec:geom})  & homography $\to$ floor distance & 26.8\% \\
MV-STGCN (\S\ref{sec:stgcn})              & graph net on projected tracks   & 25.2\% \\
Thermo-X3D (\S\ref{sec:x3d})              & 3-D CNN on raw volume           & 41.7\% \\
Image-plane voting V3 (\S\ref{sec:v3})    & box-gap + camera vote           & 28.6\% \\
\;+ box-merge / router V4 (\S\ref{sec:v4})& cleaner boxes; people routing   & 35.5\% \\
\;+ blob-merge V6.3 (\S\ref{sec:v6})      & two bodies = one warm blob      & 44.6\% \\
\;+ temporal V8 (\S\ref{sec:v8})          & morphological smoothing         & 53.8\% \\
\;+ two-body merge V9 (\S\ref{sec:v9})    & reject 3-body clusters          & 55.0\% \\
\rowcolor{rowhi}\;+ raw detector (config-D, \S\ref{sec:prep}) & raw for detection, residual for blobs & \textbf{60.8\%} \\
\midrule
\multicolumn{2}{l}{\textit{Honest cross-validated estimate of config-D (\S\ref{sec:cv})}} & $\sim$44\% \\
\bottomrule
\end{tabular}
\end{table}

% ============================================================
\section{Part A — The three designed detectors}
% ============================================================
The system design specified three contact detectors spanning a spectrum from
classical geometry to deep appearance models. We implemented and evaluated all
three first.

% ------------------------------------------------------------
\subsection{Geometric multi-view fusion}
\label{sec:geom}
\textbf{Method.} For each camera, MobileNet-SSD detects people; each box is
reduced to its \emph{foot-point} (bottom-centre $\approx$ floor contact). A
per-camera homography $H_k$ projects foot-points onto a common floor plane.
Projections are clustered and cross-camera-validated into unique
\emph{actor positions}; any pair within a distance $\delta$ is flagged as contact.

\textbf{The homography problem.} Standard calibration needs heated targets at
\emph{known} floor coordinates, which were never surveyed. We therefore
\emph{self-calibrated}: in the calibration recording one person is visible to all
three cameras, so that person's foot-point in each view is an image of one common
floor point. Across the $\sim$146-frame walk, this yields image-to-image
correspondences from which $H_{1\to0}, H_{2\to0}$ are solved by DLT inside a
RANSAC loop (camera~0 as the reference plane). Distances are then in camera-0
floor-pixels, not metres, so $\delta$ is tuned empirically.

\textbf{Result.} F1 \textbf{26.8\%} (precision 26.8\%, recall 26.8\%, FAR 4.6\%).

\confmat{Contact}{None}{11}{30}{30}{624}

\textbf{Why it underperforms.} A diagnostic on true-contact frames showed that,
with strict two-camera validation, the mean number of \emph{surviving} actors was
only \textbf{0.16} --- the self-calibrated homography is accurate enough to
localise but too noisy to fuse: the same person projects to different floor
locations in different views, so cross-camera agreement (within the clustering
radius) almost never holds. Relaxing the validation produces phantom duplicate
actors instead. Two failure modes emerge that recur throughout:
\textbf{(\#1) homography noise} and \textbf{(\#2) blob-merge} --- when two people
touch, they fuse into one detection, so there is no second box to measure.

% ------------------------------------------------------------
\subsection{MV-STGCN (multi-view spatiotemporal graph CNN)}
\label{sec:stgcn}
\textbf{Method.} The same SSD+homography front-end produces actor nodes; a Kalman
tracker (constant-velocity, correcting the $\sim$50\,ms inter-camera I\textsuperscript{2}C
phase shift) adds velocities. A sliding window of 16 actor-graphs feeds an
adaptive-adjacency ST-GCN that classifies contact/no-contact. (Because our labels
are binary with no actor positions, we injected SSD+fusion actor positions into
the training events so the graph had node features.)

\textbf{Result.} F1 \textbf{25.2\%} (precision 14.8\%, recall \textbf{82.9\%},
FAR \textbf{29.8\%}).

\confmat{Contact}{None}{34}{7}{195}{459}

\textbf{Reading.} High recall but a 30\% false-alarm rate: it fires on essentially
every multi-person non-contact scene (\texttt{1person~cig}, \texttt{2ppl\_walk},
\texttt{3ppl\_dance}). It inherits failure mode \#1 directly --- mis-projected
duplicate actors land close together and look like contact.

% ------------------------------------------------------------
\subsection{Thermo-X3D (volumetric appearance network)}
\label{sec:x3d}
\textbf{Method.} A small (2+1)D factorised X3D encoder processes each camera's
raw 16-frame thermal \emph{volume}, with CBAM thermal attention to focus on hot
voxels; the three per-camera features are late-fused and classified. It uses
\textbf{no bounding boxes and no homography} --- the designed failsafe for when
contact merges two people into one blob.

\textbf{Result.} F1 \textbf{41.7\%} (precision 48.4\%, recall 36.6\%, FAR 2.4\%)
--- the best of the three designed detectors.

\confmat{Contact}{None}{15}{26}{16}{638}

\textbf{Reading.} Being detector- and homography-independent, it avoids both
failure modes and is far cleaner (FAR 2.4\%). Its misses are on
\texttt{2ppl\_fight} (0\% recall there) --- a \emph{data} problem (only $\sim$90
positive training windows), not a method flaw. This result motivated the rest of
the work: \emph{a pixel/appearance signal beats the geometric pipeline}.

\medskip\noindent\textbf{Part A takeaway.} Homography is the weak link on this
data; the best detector ignores it. We then asked whether a \emph{simpler,
homography-free} detector could do better.

% ============================================================
\section{Part B — A homography-free image-plane family}
% ============================================================
The idea: run SSD per camera, decide ``touch'' from the 2-D boxes in each image,
and combine the three views. No homography, no world projection.

% ------------------------------------------------------------
\subsection{Image-plane voting (V1/V2/V3)}
\label{sec:v3}
\textbf{Method.} A camera ``votes touch'' if two person-boxes have an
edge-to-edge gap below $\tau$ pixels (0 = overlapping). Three ways to combine the
three votes: \textbf{V1 majority} ($\ge$2 cameras), \textbf{V2 unanimous} (all 3),
\textbf{V3 relaxed} ($\ge$2 vote, and the third corroborates --- near-touch or a
single merged blob --- otherwise a third camera clearly seeing two separated
people vetoes). $\tau$ tuned on validation.

\textbf{Result.}
\begin{center}\renewcommand{\arraystretch}{1.15}
\begin{tabular}{l c c c c}
\toprule
& \textbf{Prec} & \textbf{Rec} & \textbf{F1} & \textbf{FAR} \\
\midrule
V1 majority   & 18.5\% & 56.1\% & 27.9\% & 15.4\% \\
V2 unanimous  & 14.7\% & 53.7\% & 23.0\% & 19.6\% \\
\textbf{V3 relaxed} & 19.2\% & 56.1\% & \textbf{28.6\%} & 14.8\% \\
\bottomrule
\end{tabular}\end{center}

\textbf{Reading.} V3 is best. Counter-intuitively V2 (unanimous) is \emph{worst}:
requiring strict 3-camera agreement is so rare that tuning loosens $\tau$ to
compensate, which floods false positives. Already comparable to the geometric
detector, with no homography.

% ------------------------------------------------------------
\subsection{Error analysis of V3}
\label{sec:erranalysis}
We dissected V3's 18 missed contacts (FN) and 97 false alarms (FP):
\begin{itemize}[leftmargin=1.4em, itemsep=1pt]
\item \textbf{FN cause 1 --- blob-merge} (11 frames, all \texttt{2ppl\_fight}):
      the fighters fuse into one box in 2 of 3 cameras (only one camera still
      sees two), so V3 gets one TOUCH vote and needs two.
\item \textbf{FN cause 2 --- threshold strictness} (7 frames,
      \texttt{3pp\_surprise}): boxes are 1--6\,px apart, just above $\tau$.
\item \textbf{FP cause 1 --- crowd proximity} ($\sim$86 FP: \texttt{3pplhedroncolider},
      \texttt{3ppl\_dance}, \texttt{the\_more\_the\_merrier}): three clustered
      people whose boxes overlap in every view without contact.
\item \textbf{FP cause 2 --- detector over-segmentation} ($\sim$7 FP,
      \texttt{1person~cig}): one person (or person + cigarette) split into two
      abutting boxes.
\end{itemize}
These four causes drove every subsequent improvement.

% ------------------------------------------------------------
\subsection{Targeted improvements V3.1--V3.5}
\label{sec:v3x}
Each idea was implemented and evaluated \emph{separately} (re-tuned on validation):

\begin{center}\renewcommand{\arraystretch}{1.15}
\begin{tabular}{l l c c c}
\toprule
& \textbf{Idea} & \textbf{Rec} & \textbf{F1} & \textbf{FAR} \\
\midrule
V3.0 & baseline                         & 56.1\% & 28.6\% & 14.8\% \\
V3.1 & merge-aware quorum               & \textbf{92.7\%} & 28.3\% & 29.1\% \\
V3.2 & temporal persistence ($\ge$2/3)  & 43.9\% & 25.9\% & 12.2\% \\
V3.3 & size-normalised gap              & 75.6\% & 30.2\% & 20.3\% \\
\textbf{V3.4} & merge over-segmented boxes & 58.5\% & \textbf{32.0\%} & 13.0\% \\
V3.5 & OR Thermo-X3D ensemble           & 36.6\% & 42.3\% & 2.3\% \\
\bottomrule
\end{tabular}\end{center}

\textbf{Reading.} \textbf{V3.4 (box-merge)} is the only free win: merging
over-segmented boxes raised both precision and recall (F1 +3.4) by removing
single-person false alarms. \textbf{V3.1 (merge-aware)} --- counting a single
merged blob toward the quorum --- nearly eliminates the fight misses (recall
92.7\%) but at a heavy precision cost. \textbf{V3.2 (persistence)} hurt (contacts
are short). \textbf{V3.5} degenerated to Thermo-X3D alone: naively OR-ing the box
detector with X3D adds only false positives, so tuning zeroed out the box branch.

% ------------------------------------------------------------
\subsection{Compound V4 / V5: merge-aware + box-merge + people-count routing}
\label{sec:v4}
\textbf{Method.} Combine the winners: clean boxes (V3.4), merge-aware quorum
(V3.1), and route by crowd size --- trust the box logic on $\le$2-person frames
(where it catches fights and X3D fails), defer to Thermo-X3D on $\ge$3-person
frames (where the box logic sprays). V5 is the same routing expressed as an
explicit people-count switch.

\textbf{Result.} F1 \textbf{35.5\%} (precision 24.3\%, recall 65.9\%, FAR 12.8\%);
V5 is numerically identical to V4. The crowd-veto eliminated \texttt{3ppl\_dance}
false alarms, and merge-awareness recovered \texttt{2ppl\_fight} (84\% recall ---
the scene X3D misses entirely). Remaining false alarms shifted to single/pair
scenes where SSD detects a hot \emph{object} (cigarette, heater) as a second
``person'' --- which pure box geometry cannot distinguish from two people.

% ------------------------------------------------------------
\subsection{V6 --- putting the thermal signal back in}
\label{sec:v6}
The box family ignored temperature; both its remaining error types are thermal
questions. We tested four thermal cues separately:

\begin{center}\renewcommand{\arraystretch}{1.15}
\begin{tabular}{l l c c c}
\toprule
& \textbf{Idea} & \textbf{Rec} & \textbf{F1} & \textbf{FAR} \\
\midrule
V6.1 & heat-bridge (warm path between boxes)        & 36.6\% & 39.0\% & 3.2\% \\
V6.2 & temperature box-rejection (drop hot objects) & 65.9\% & 41.2\% & 9.6\% \\
\rowcolor{rowhi}\textbf{V6.3} & \textbf{blob-merge (two people in one warm component)} & 61.0\% & \textbf{44.6\%} & 7.0\% \\
V6.4 & learned logistic combiner                    & 41.5\% & 35.1\% & 6.0\% \\
\bottomrule
\end{tabular}\end{center}

\textbf{Method \& reading.} \textbf{V6.3} is the breakthrough and the first
homography-free detector to \emph{beat} Thermo-X3D. It segments warm connected
components on the background-subtracted residual and flags contact when two
detected people fall in the \emph{same} component (their bodies have fused). This
single idea fixes both error types at once: a cigarette is a \emph{separate}
component from its smoker (object false alarms collapse, \texttt{1person~cig}
23$\to$4 FP), and people merely standing close are separated by cool air
(\texttt{3ppl\_dance} $\to$0 FP). \textbf{V6.1 (heat-bridge)} sampling temperature
along the line between two boxes gave the best precision/FAR but cut recall (too
strict). \textbf{V6.4 (learned combiner)} \emph{underperformed} the hand rule ---
the first sign that learning loses on this little data.

% ------------------------------------------------------------
\subsection{V7 --- blob-merge on crowds}
\label{sec:v7}
\textbf{Method.} V6.3 deferred 3-person crowds to X3D, which fails
\texttt{3pplhedroncolider}. V7 applies the blob-merge test to all regimes.

\textbf{Result.} F1 \textbf{44.6\%} (V7b, recall 65.9\%). It recovered
\texttt{3pplhedroncolider} recall from 0\% to 86\%, but that scene's
\emph{negative} frames also trip the test (three colliding bodies are one blob),
so the gained true positives were offset by new false positives --- F1 unchanged.
V7b is preferred over V6.3 at equal F1 because it covers all four contact scenes,
including collisions.

% ------------------------------------------------------------
\subsection{V8 --- temporal processing}
\label{sec:v8}
\textbf{Method.} Contacts persist over time, so the per-frame errors are largely
temporal noise. We applied, on the per-scene prediction stream:
\textbf{(T1) morphological open+close} (remove short positive runs = flicker FPs;
fill short negative gaps = dropout FNs); \textbf{(T2) centred window vote};
\textbf{(T3) a temporal logistic regression} on windowed features.

\begin{center}\renewcommand{\arraystretch}{1.15}
\begin{tabular}{l c c c}
\toprule
& \textbf{Rec} & \textbf{F1} & \textbf{FAR} \\
\midrule
\textbf{T1 morphology open+close} & 68.3\% & \textbf{53.8\%} & 5.4\% \\
T2 centred window vote            & 68.3\% & 51.4\% & 6.1\% \\
T3 temporal logistic regression   & 36.6\% & 34.9\% & 4.6\% \\
\bottomrule
\end{tabular}\end{center}

\textbf{Reading.} \textbf{T1} lifted F1 by \textbf{+9 points} (44.6$\to$53.8),
improving precision \emph{and} recall. T3 (learned) again lost --- the second
confirmation that learned additions overfit here.

% ------------------------------------------------------------
\subsection{V9 --- two-body merge + temporal}
\label{sec:v9}
\textbf{Method.} The remaining gap was \texttt{3pplhedroncolider}'s sustained
false positives. V9 adds a spatial discriminator: on crowds, fire only when a
warm component contains \emph{exactly two} person-centres (a genuine pairwise
merge), rejecting the three-body cluster; then T1 temporal smoothing.

\textbf{Result.} F1 \textbf{55.0\%} (precision 56.4\%, recall 53.7\%, FAR 2.6\%).
\texttt{3pplhedroncolider} false alarms dropped 19$\to$4. A
confidence-aware variant trades to recall 73.2\% at F1 54.1\%.

% ============================================================
\section{Part C — Preprocessing ablation (key finding)}
\label{sec:prep}
% ============================================================
\textbf{Question.} The Tateno preprocessing (Gaussian smooth $\to$ background
subtract $\to$ residual) feeds the pipeline in two places: the SSD detector and
the blob segmentation. Does it help? We retrained the SSD on \emph{raw} frames
and tested all combinations with the V9 detector (no X3D):

\begin{center}\renewcommand{\arraystretch}{1.15}
\begin{tabular}{l l c}
\toprule
\textbf{Config} & \textbf{Detector / blob input} & \textbf{F1} \\
\midrule
A & Tateno SSD\, /\, residual blob & 54.5\% \\
B & Tateno SSD\, /\, raw-temp blob & 8.7\% \\
C & raw SSD\, /\, raw-temp blob    & 0.0\% \\
\rowcolor{rowhi}\textbf{D} & \textbf{raw SSD\, /\, residual blob} & \textbf{60.8\%} \\
\bottomrule
\end{tabular}\end{center}

\textbf{Finding --- preprocessing should differ by task.}
\begin{itemize}[leftmargin=1.4em, itemsep=1pt]
\item The blob test \emph{needs} the residual (B, C collapse to $\sim$0): on raw
      temperature, bodies barely exceed the warm background, so connected-component
      segmentation cannot separate people.
\item The residual \emph{hurts} the detector (D 60.8\% $>$ A 54.5\%): subtracting
      the background removes body signature the detector uses; the raw-trained SSD
      even reached a lower training loss.
\end{itemize}
\textbf{config-D} --- raw SSD + residual blob + two-body merge + temporal --- is
the best detector at \textbf{F1 60.8\%} (precision 50.8\%, recall 75.6\%,
FAR 4.6\%) on this split, with no homography and no contact training labels.

% ============================================================
\section{Part D — Reaching 65\%? The honest reality check}
\label{sec:cv}
% ============================================================

\subsection{Joint re-tuning overfits (V10)}
Re-tuning all parameters jointly (score, blob threshold, $\tau_2$, X3D gate,
morphology) found a validation optimum (val F1 55.4\%) that \emph{dropped to
54.4\% on test} --- below config-D's 60.8\%. With only 40 validation positives,
tuning $\sim$6 knobs fits noise.

\subsection{Cross-validation: the true number}
We ran 4-fold timeline cross-validation over the 1108-frame val+test pool (81
positives, every frame tested once, no detector leakage):

\begin{center}\renewcommand{\arraystretch}{1.15}
\begin{tabular}{l c c c}
\toprule
\textbf{Detector} & \textbf{fold F1s} & \textbf{mean $\pm$ std} & \textbf{pooled F1} \\
\midrule
config-D rule        & 32/51/65/0\,\% & 36.9\% $\pm$ 24.3\% & \textbf{44.3\%} \\
logistic (no motion) & 34/58/42/0\,\% & 33.4\% $\pm$ 21.1\% & 42.6\% \\
logistic (+ motion)  & 29/52/36/0\,\% & 29.4\% $\pm$ 19.0\% & 37.8\% \\
\bottomrule
\end{tabular}\end{center}

\textbf{Findings.} (i) config-D's honest F1 is $\approx$\textbf{44\%}, not 60.8\%
--- the single split was a favourable region. (ii) The fold-to-fold spread is
enormous (\textbf{$\pm$24 points}, folds ranging 0--65\%). (iii) \textbf{Motion
features hurt} ($-4$ points) --- frame-difference energy and merge run-length
added noise the model overfit. (iv) Every learned variant again trailed the rule.

\subsection{How much data would 65\% actually need?}
A bootstrap of config-D's pooled predictions gives 95\% confidence intervals:
frame-level \textbf{[35.7\%, 52.4\%]} ($\pm$8.3) and scene-level
\textbf{[12.2\%, 66.3\%]} ($\pm$27). The scene-level interval --- spanning almost
the whole range because positives come from only four scenes --- is the binding
constraint. To estimate F1 to $\pm$2.5\% (so 65\% is distinguishable from 60\%)
needs on the order of \textbf{$\sim$900 positive frames ($\sim$11$\times$ current)}
and \textbf{3--4$\times$ more distinct contact scenes}.

\subsection{Can old MLX90640 touch data substitute?}
We pretrained Thermo-X3D on the older MLX90640 touch recordings (real contact, but
24$\times$32, upsampled to 62$\times$80) and fine-tuned on Waveshare. It
\textbf{reduced} F1 by 6.2 points (scratch 35.9\% $\to$ transfer 29.7\%): the low
resolution cannot reproduce the two-body structure the method relies on, and the
sensor-domain gap caused over-firing. Synthetic compositing of touches was judged
unviable (it must be geometrically consistent across all three cameras
simultaneously).

% ============================================================
\section{Conclusions}
% ============================================================
\begin{enumerate}[leftmargin=1.6em, itemsep=2pt]
\item \textbf{A practical homography-free detector exists.} ``config-D''
      (MobileNet-SSD + residual blob-merge + two-body test + temporal smoothing)
      reaches F1 60.8\% on a matched split and beats the deep Thermo-X3D baseline,
      while being explainable and needing no contact training labels.
\item \textbf{Transferable design lessons:} (a) the connected-warm-component
      \emph{merge} test is the single most informative contact signal; (b)
      preprocessing should differ by task --- \emph{raw} frames for the detector,
      \emph{residual} for blob segmentation; (c) temporal smoothing adds $\sim$9
      F1 points; (d) homography self-calibration is too noisy to rely on here.
\item \textbf{Robust negative results:} every learned addition (logistic
      combiner, motion features, MLX transfer) and aggressive joint tuning matched
      or underperformed the simple rule --- all symptoms of too few examples.
\item \textbf{The system is data-limited, not algorithm-limited.} Cross-validation
      gives F1 $\approx$44\% with a [12\%, 66\%] confidence interval; 188 positives
      in four scenes cannot reliably train richer models or measure performance.
\item \textbf{Recommendation.} Adopt config-D as the working method and prioritise
      \textbf{collecting more contact recordings across more distinct scenarios}
      (target several times the current positive count and contact scenes). This is
      the single change expected to both raise and stabilise performance; further
      tuning on the present data is within measurement noise.
\end{enumerate}

\end{document}
